\chapter{Moteur de Similarité et Recherche par l'Image}
\label{chap:recherche}

\section{Espaces Métriques et Mesures de Proximité}

La recherche par le contenu nécessite de projeter les images au sein d'un espace vectoriel normé et d'évaluer leur proximité relative. Cette fonction est assurée par le composant \texttt{SimilarityEngine} (\texttt{src/services/similarity\_engine.py}).

Soit $\mathbf{q} \in \mathbb{R}^D$ le vecteur représentant la requête visuelle de l'utilisateur et $\mathbf{c} \in \mathbb{R}^D$ le vecteur correspondant à une image candidate issue de la collection indexée.

Toutes les mesures de proximité dans Perceptra sont calibrées pour appartenir strictement à l'intervalle $[0, 1]$, où la valeur $1.0$ indique une identité visuelle parfaite et $0.0$ une divergence totale.

\begin{infobox}{Métrique Cosinus Normalisée (Défaut)}
La similarité cosinus mesure l'angle géométrique entre deux vecteurs indépendamment de leur magnitude globale, ce qui s'avère particulièrement pertinent pour comparer des histogrammes de fréquences ou de textures~:
\begin{equation}
\cos(\mathbf{q}, \mathbf{c}) = \frac{\mathbf{q} \cdot \mathbf{c}}{\|\mathbf{q}\|_2 \cdot \|\mathbf{c}\|_2} \implies S_{\text{cos}}(\mathbf{q}, \mathbf{c}) = \frac{1 + \cos(\mathbf{q}, \mathbf{c})}{2} \in [0, 1]
\end{equation}
Les évaluations expérimentales sur GHIM-10K confirment la supériorité systématique de la similarité cosinus, surpassant la distance euclidienne de plus de 10 points de précision ($P@5 = 52,5\,\%$ contre $42,0\,\%$).
\end{infobox}

\begin{infobox}{Métrique Euclidienne Inversée}
Lorsque l'utilisateur sélectionne la métrique euclidienne $L_2$, la distance géométrique $d_2(\mathbf{q}, \mathbf{c}) = \|\mathbf{q} - \mathbf{c}\|_2$ est convertie en score de similarité borné selon la décroissance harmonique standard~:
\begin{equation}
S_{\text{euc}}(\mathbf{q}, \mathbf{c}) = \frac{1}{1 + \|\mathbf{q} - \mathbf{c}\|_2} \in [0, 1]
\end{equation}
\end{infobox}

\section{Agrégation Convexe Multi-Descripteurs}

Pour éviter l'écueil classique de la «~malédiction de la dimensionnalité~» induit par la concaténation brute d'espaces hétérogènes, Perceptra maintient chaque descripteur dans son sous-espace métrique propre.

Soit $\mathcal{F}$ l'ensemble des familles de descripteurs actives. Pour chaque descripteur $j \in \mathcal{F}$, nous disposons d'un vecteur de requête $\mathbf{q}_j$, d'un candidat $\mathbf{c}_j$ et d'un poids $w_j \ge 0$. Les poids actifs sont préalablement normalisés sous contrainte convexe~:
\begin{equation}
\tilde{w}_j = \frac{w_j}{\sum_{k \in \mathcal{F}} w_k}, \quad \text{tel que } \sum_{j \in \mathcal{F}} \tilde{w}_j = 1.0
\end{equation}

La similarité globale $S(\mathbf{q}, \mathbf{c})$ entre l'image requête et l'image candidate est alors formulée comme la combinaison linéaire convexe des similarités individuelles~:
\begin{equation}
S(\mathbf{q}, \mathbf{c}) = \sum_{j \in \mathcal{F}} \tilde{w}_j \cdot S_j(\mathbf{q}_j, \mathbf{c}_j)
\end{equation}

\section{Décomposition Analytique et Ordonnancement Top-K}

Pour offrir une transparence maximale à l'utilisateur, la méthode \texttt{aggregate\_multi\_feature\_similarity} extrait simultanément la décomposition des affinités par famille conceptuelle~:
\begin{itemize}
    \item \textbf{Similarité Chromatique ($S_{\text{color}}$)~:} Moyenne arithmétique des scores des composantes de couleur actives (\texttt{color\_histogram}, \texttt{color\_moments}, \texttt{dominant\_colors})~;
    \item \textbf{Similarité Texturale ($S_{\text{texture}}$)~:} Moyenne arithmétique des scores des composantes texturales actives (\texttt{glcm}, \texttt{lbp\_histogram})~;
    \item \textbf{Scores fins par descripteur~:} Dictionnaire associatif $\{j: S_j(\mathbf{q}_j, \mathbf{c}_j)\}$.
\end{itemize}

Le service de requête \texttt{QueryService} (\texttt{src/services/query\_service.py}) orchestre l'exécution de la recherche en mémoire~: filtrage du périmètre (\texttt{all}, \texttt{ghim10k}, \texttt{my\_images}), balayage matriciel vectorisé sous NumPy évaluant 10\,000 vecteurs en moins de 45 millisecondes, tri décroissant et extraction des $K$ meilleurs résultats (\texttt{Top-K}).
